\begin{table}[htbp]
\centering
\small
\begin{threeparttable}
\caption{Sentence length and Flesch reading ease by predicted CEFR level}
\label{tab:readability-by-level}
\begin{tabular}{llrrr}
\hline
Language & Level & n & Sentence length & Flesch \\
\hline
German & A & 1,105 & 13.0 [10.0, 17.0] & 24.7 [9.5, 37.7] \\
 & B1 & 107 & 12.0 [9.8, 15.4] & 35.0 [22.6, 47.9] \\
 & B2 & 280 & 12.7 [11.1, 15.0] & 40.5 [29.4, 46.3] \\
 & C & 3,507 & 14.3 [12.5, 17.3] & 36.1 [25.4, 43.0] \\
\hline
Dutch & A & 518 & 9.4 [1.0, 11.3] & 47.3 [-25.1, 65.3] \\
 & B1 & 1,043 & 11.6 [10.5, 13.0] & 53.7 [46.4, 62.9] \\
 & B2 & 772 & 16.0 [13.1, 18.0] & 42.8 [36.8, 48.8] \\
 & C & 77 & 18.5 [16.9, 20.2] & 37.9 [28.0, 44.9] \\
\hline
English & A & 204 & 14.5 [12.3, 17.5] & 39.3 [29.5, 46.5] \\
 & B1 & 505 & 18.2 [15.7, 20.5] & 37.5 [28.4, 47.8] \\
 & B2 & 734 & 20.9 [18.4, 23.5] & 35.1 [26.9, 42.4] \\
 & C & 145 & 24.3 [22.2, 28.5] & 24.3 [17.1, 33.3] \\
\hline
\end{tabular}
\begin{tablenotes}
\footnotesize
\item Medians with interquartile range in brackets. Flesch variants: Amstad (German), Douma (Dutch), Flesch (English); higher means easier. Syllables use a vowel-cluster heuristic.
\item Spearman correlation of predicted level (1 to 4) with sentence length and Flesch score: German 0.19 and 0.18; Dutch 0.66 and -0.23; English 0.53 and -0.23.
\end{tablenotes}
\end{threeparttable}
\end{table}
